\documentclass[12pt]{article}
\usepackage[utf8]{inputenc}
\usepackage[T1]{fontenc}
\usepackage[margin=1in]{geometry}
\usepackage{setspace}
\usepackage{xcolor}
\usepackage[colorlinks=true, urlcolor=blue!60!black, linkcolor=blue!60!black]{hyperref}
\setlength{\parindent}{2em}
\setlength{\parskip}{1.25em}
\renewcommand{\baselinestretch}{1.0}
\title{\textit{Ward v. Rock Against Racism}, 491 U.S. 781 (1989)}
\author{Case Brief}
\date{}
\begin{document}
\maketitle

\section{Case Metadata}
\begin{description}
  \item[Full Case Name:] \textit{Benjamin Ward, Police Commissioner of the City of New York, et al. v. Rock Against Racism}
  \item[Citation:] 491 U.S. 781 (1989) (Docket No. 88-226)
  \item[Argued / Decided:] February 27, 1989 / June 22, 1989
  \item[Vote:] 6--3 (5--1--3)
  \item[Author:] Justice Anthony M. Kennedy (joined by Chief Justice Rehnquist and Justices White, O'Connor, and Scalia)
  \item[Concurrences / Dissents:] Justice Harry A. Blackmun (concurring in the result, without opinion); Justice Thurgood Marshall (dissenting, joined by Justices Brennan and Stevens)
  \item[Chapter / Topic:] Free Speech \& Press --- Content-Neutral Time, Place, and Manner Regulations in Traditional Public Forums
  \item[Primary Source:] \href{https://supreme.justia.com/cases/federal/us/491/781/}{Justia U.S. Supreme Court Center (\textit{Ward v. Rock Against Racism}, 491 U.S. 781)}
\end{description}

\section{Facts and Procedural History}
Rock Against Racism (RAR), an unincorporated association dedicated to promoting antiracist views, sponsored annual programs of speeches and rock music from 1979 through 1986 at the Naumburg Acoustic Bandshell (spelled ``Naumberg'' in the Court's opinion) in New York City's Central Park. The bandshell faces west toward the Sheep Meadow---a city-designated quiet zone for passive recreation---and nearby residential apartments along Central Park West. Over several years, RAR's concerts generated numerous excessive-noise complaints from park users and nearby residents. At RAR's 1984 concert, performers ignored repeated warnings and two citations for excessive volume; when police ultimately shut off power to the sound system, the crowd became abusive and disruptive. Meanwhile, other bandshell sponsors using inadequate equipment or unskilled technicians produced insufficient sound levels, leading to disappointed or unruly crowds and prompting some performers to raise volume to compensate for poor sound mix.

After rejecting fixed decibel caps (because sound propagation varies with temperature, foliage, and crowd size) and rejecting having a city technician operate unfamiliar sponsor-supplied equipment, New York City promulgated Use Guidelines on March 21, 1986. The sound-amplification guideline required all bandshell performers to use high-quality amplification equipment leased by the Parks Department and operated by an independent, experienced sound technician retained by the city. RAR sued in the U.S. District Court for the Southern District of New York, challenging the guideline as facially invalid under the First Amendment. After a five-day trial, the District Court upheld the guideline, 658 F. Supp. 1346 (S.D.N.Y. 1987), finding that the city's system met RAR's technical needs, that city policy gave sponsors autonomy over sound ``mix'' (balancing instruments and vocals), and that city officials conferred with sponsors before lowering volume. The Second Circuit reversed, 848 F.2d 367 (2d Cir. 1988), holding that time, place, and manner regulations must be the ``least intrusive'' means reasonably necessary to achieve the city's goal and pointing to less restrictive alternatives such as volume-limiting devices or pulling the plug. The Supreme Court granted certiorari and reversed.

\section{Issue(s) Presented}
\begin{enumerate}
  \item Is New York City's Central Park bandshell sound-amplification guideline---requiring performers to use city-provided sound equipment and a city-retained technician---a content-neutral regulation of the time, place, and manner of protected speech?
  \item Does the ``narrow tailoring'' requirement for content-neutral time, place, and manner regulations in a traditional public forum require the government to adopt the least restrictive or least intrusive means of achieving its substantial interest?
\end{enumerate}

\section{Holding}
Yes to Question 1; No to Question 2. Reversing the Second Circuit, the Supreme Court held that music is protected expression under the First Amendment and that the bandshell is a traditional public forum, but upheld the city's sound-amplification guideline as a valid, content-neutral time, place, and manner regulation. Crucially, the Court clarified that narrow tailoring under intermediate scrutiny does \textit{not} require the ``least restrictive or least intrusive means''; rather, a regulation is narrowly tailored so long as it promotes a substantial government interest that would be achieved less effectively absent the regulation, and the means chosen do not burden substantially more speech than necessary.

\section{Reasoning}
Writing for a five-Justice majority, Justice Kennedy began by affirming that ``[m]usic, as a form of expression and communication, is protected under the First Amendment,'' and treated the bandshell as a public forum governed by the three-part time, place, and manner framework from \textit{Clark v. Community for Creative Non-Violence}, 468 U.S. 288, 293 (1984).

\textbf{1. Content Neutrality and Unbridled Discretion:} The Court stated that the ``principal inquiry in determining content-neutrality'' is whether the government adopted a regulation ``because of disagreement with the message it conveys,'' making the ``government's purpose'' the controlling consideration. New York's primary objective---controlling excessive noise to preserve the tranquility of the Sheep Meadow and adjacent residences---had nothing to do with expressive content. Nor did the city's secondary interest in ensuring adequate sound quality make the rule content-based: the city expressly disavowed imposing aesthetic standards and required its technician to defer to sponsors on sound mix. Rejecting RAR's facial challenge that the guideline vested unbridled discretion in city officials (unlike licensing schemes that allow officials to deny permits altogether, as in \textit{City of Lakewood v. Plain Dealer Publishing Co.}, 486 U.S. 750 (1988)), the Court held that the guideline's text and the city's narrowing administrative construction adequately constrained official discretion and did not constitute a prior restraint.

\textbf{2. Narrow Tailoring Without Least-Restrictive-Means Review:} The city possessed substantial interests in protecting citizens from unwelcome noise---even in public parks---and in ensuring sufficient amplification so concertgoers could hear performances. Faulting the Second Circuit for importing a ``least intrusive means'' test from strict scrutiny (or misreading \textit{United States v. O'Brien}, 391 U.S. 367 (1968)), Justice Kennedy reaffirmed that least-restrictive-alternative analysis ``has never been a part of the inquiry into the validity of a time, place, and manner regulation.'' Instead, narrow tailoring is satisfied if the regulation promotes a substantial interest that ``would be achieved less effectively absent the regulation,'' provided the government does not ``burden substantially more speech than is necessary.'' Courts must defer to the responsible decisionmaker's reasonable choice among effective regulatory methods and evaluate the rule's fit against the ``overall problem'' rather than a single speaker's circumstances. Because the city's technician deferred to performers on sound mix, the guideline caused no material interference with artistic expression.

\textbf{3. Ample Alternative Channels:} Finally, the guideline easily left open ample alternative channels of communication: it did not ban any type of music or expression at the bandshell, and merely moderating maximum volume did not render remaining avenues of expression inadequate.

\section{Concurring and Dissenting Opinions}
\begin{itemize}
  \item \textbf{Justice Blackmun (Concurring in the Result):} Concurred in the judgment reversing the Second Circuit without issuing a separate opinion.
  \item \textbf{Justice Marshall (Dissenting, joined by Justices Brennan and Stevens):} Agreed that the guideline was content-neutral and served a significant interest in controlling noise, but forcefully condemned the majority for ``eviscerat[ing]'' narrow tailoring and replacing constitutional scrutiny with ``mandatory deference.'' Citing \textit{Schneider v. State}, 308 U.S. 147 (1939), and \textit{Martin v. Struthers}, 319 U.S. 141 (1943), Justice Marshall argued that narrow tailoring historically required weighing a regulation's incremental efficacy against its extra burden on speech when less intrusive tools---such as enforcing existing noise ordinances or perimeter decibel monitoring, as used by the National Park Service---were available. Under the majority's rule, any regulation that advances a goal ``in the slightest'' survives unless it is gratuitous. Second, Justice Marshall characterized exclusive municipal control over sound equipment as a classic prior restraint---analogous to Crown monopolies over printing presses, effecting censorship ``by a single turn of a knob'' without definite standards or prompt judicial review under \textit{Freedman v. Maryland}, 380 U.S. 51 (1965).
\end{itemize}

\section{Key Quotations}
\begin{quote}
``Music is one of the oldest forms of human expression. \dots\ Music, as a form of expression and communication, is protected under the First Amendment.'' \\
--- Justice Kennedy, for the Court, \href{https://supreme.justia.com/cases/federal/us/491/781/}{491 U.S. at 790}.
\end{quote}

\begin{quote}
``The principal inquiry in determining content-neutrality, in speech cases generally and in time, place, or manner cases in particular, is whether the government has adopted a regulation of speech because of disagreement with the message it conveys. The government's purpose is the controlling consideration.'' \\
--- Justice Kennedy, for the Court, \href{https://supreme.justia.com/cases/federal/us/491/781/}{491 U.S. at 791}.
\end{quote}

\begin{quote}
``Lest any confusion on the point remain, we reaffirm today that a regulation of the time, place, or manner of protected speech must be narrowly tailored to serve the government's legitimate, content-neutral interests, but that it need not be the least restrictive or least intrusive means of doing so.'' \\
--- Justice Kennedy, for the Court, \href{https://supreme.justia.com/cases/federal/us/491/781/}{491 U.S. at 798--799}.
\end{quote}

\begin{quote}
``Under more familiar prior restraints, government officials censor speech `by a simple stroke of the pen' \dots\ Here, it is done by a single turn of a knob. \dots\ Today, the majority enshrines efficacy, but sacrifices free speech.'' \\
--- Justice Marshall, dissenting, \href{https://supreme.justia.com/cases/federal/us/491/781/}{491 U.S. at 807--808}.
\end{quote}

\section{Doctrinal Significance and Modern Legacy (Through 2026)}
\textit{Ward v. Rock Against Racism} supplies the canonical modern formula for intermediate scrutiny of content-neutral speech regulations in public forums (and, via \textit{Turner Broadcasting System, Inc. v. FCC}, 512 U.S. 622 (1994), for content-neutral media regulations generally). Over the past decade, however, the Supreme Court has refined both halves of \textit{Ward}:
\begin{itemize}
  \item \textbf{Restoring Bite to Narrow Tailoring (\textit{McCullen v. Coakley}, 573 U.S. 464 (2014)):} Applying \textit{Ward}, a unanimous Court struck down Massachusetts's 35-foot buffer zone around abortion clinics. While reiterating that \textit{Ward} does not demand the least restrictive means, Chief Justice Roberts emphasized that the government ``may not sacrifice speech for efficiency'' and must show it ``seriously undertook to address the problem with less intrusive tools readily available to it'' (such as targeted injunctions or anti-obstruction laws).
  \item \textbf{Refining Content Neutrality (\textit{Reed v. Town of Gilbert}, 576 U.S. 155 (2015), and \textit{City of Austin v. Reagan National Advertising of Austin, LLC}, 596 U.S. 61 (2022)):} \textit{Reed} curtailed \textit{Ward}'s statement that ``purpose is the controlling consideration,'' holding that \textit{facially} content-based laws trigger strict scrutiny regardless of benign governmental purpose, while keeping \textit{Ward}'s purpose inquiry as a second, independent trigger for facially neutral laws. In \href{https://supreme.justia.com/cases/federal/us/596/20-1029/}{\textit{City of Austin}}, the Court tempered \textit{Reed}, holding 6--3 that an on-/off-premises sign distinction is facially content-neutral and remanding for intermediate scrutiny under \textit{Ward}.
  \item \textbf{2025--2026 Applications:} In \href{https://www.law.cornell.edu/supremecourt/text/24-656}{\textit{TikTok Inc. v. Garland}}, 604 U.S. \_\_\_ (Jan. 17, 2025), the Court cited \textit{Ward} (491 U.S. at 791) in holding that the federal divest-or-ban statute was facially content-neutral and justified by a content-neutral national-security rationale, upholding it under intermediate scrutiny. Remarkably, in \href{https://www.law.cornell.edu/supremecourt/text/24-43}{\textit{West Virginia v. B. P. J.}}, 609 U.S. \_\_\_ (June 30, 2026), the Court even imported \textit{Ward}'s aggregate-fit principle (491 U.S. at 801) into Equal Protection intermediate scrutiny, quoting \textit{Ward} for the rule that a regulation's validity depends on its relation to the ``overall problem'' rather than an individual application.
\end{itemize}

\section{Pedagogical Critiques and Exam Hypotheticals}
\textbf{Pedagogical Critiques:}
\begin{itemize}
  \item \textit{The Tension Between \textit{Ward} and \textit{McCullen}:} Students should compare \textit{Ward}'s deference to New York's rejection of decibel caps with \textit{McCullen}'s demand that Massachusetts prove less intrusive enforcement tools failed. Does \textit{McCullen} quietly revive Justice Marshall's balancing inquiry whenever a regulation imposes a heavy practical burden on core one-on-one political or sidewalk speech?
  \item \textit{Reliance on Unwritten Narrowing Constructions:} \textit{Ward} upheld broad regulatory text (``best sound'') partly because city officials testified they deferred to performers on sound mix. Critics note this weakens the facial-challenge safeguard against unbridled discretion when administrations change.
\end{itemize}

\textbf{Exam Hypothetical (and ``Opposite Side'' Contrast):} \\
\textit{Hypothetical:} Following loud late-night protests outside municipal buildings, City X enacts an ordinance governing its downtown Public Plaza (a traditional public forum). The ordinance bans all privately owned bullhorns, amplifiers, and speakers after 7:00 p.m., requiring any evening speaker seeking amplification to use a city-operated PA system capped at 70 decibels and staffed by a municipal technician who may mute the microphone if sound ``unreasonably disturbs nearby commerce or residential tranquility.'' A grassroots labor group challenging the ordinance argues it violates the First Amendment.

\textit{Analysis \& ``Opposite Side'' Contrast:}
\begin{itemize}
  \item \textit{City's Defense (Under \textit{Ward}):} City X will argue the rule is facially content-neutral and justified solely by noise control; that a 70-decibel city PA system directly advances its substantial interest in residential and commercial tranquility more effectively than after-the-fact citations; that under \textit{Ward}, it need not adopt the least restrictive means; and that ample alternative channels remain because speakers may still protest unamplified, use the city PA system, or use private amplification before 7:00 p.m.
  \item \textbf{Opposite Side Contrast (Challenger's Winning Angles Under \textit{McCullen} and Marshall's Dissent):} Unlike \textit{Ward}---where the city's system met all performers' technical needs, covered a fixed bandshell, and carried an established narrowing construction requiring deference to speakers---City X's ordinance grants a technician standardless authority to ``mute'' live political speech based on subjective ``disturbance'' with no narrowing construction or consultation protocol, inviting a strong facial unbridled-discretion/prior-restraint attack. Moreover, under \textit{McCullen v. Coakley}, if City X never tried enforcing its general noise ordinance or issuing citations to actual violators before banning all private bullhorns across an entire public plaza, the ordinance burdens substantially more peaceful speech than necessary and fails \textit{Ward}'s narrow-tailoring prong.
\end{itemize}

\end{document}
